\documentclass[11pt]{article}

% --- packages (dependency-light; compiles with tectonic or a standard TeX install) ---
\usepackage[utf8]{inputenc}
\usepackage[T1]{fontenc}
\usepackage{lmodern}
\usepackage{microtype}
\usepackage{amsmath,amssymb}
\usepackage{graphicx}
\usepackage{booktabs}
\usepackage[margin=1in]{geometry}
\usepackage{caption}
\usepackage[hidelinks]{hyperref}

\graphicspath{{./}{figures/}{../figures/}}   % ./ so a flat arXiv upload (figures beside paper.tex) also builds
\emergencystretch=3em   % absorb the few tight lines with unbreakable inline-math tuples

\newcommand{\Is}{I^{*}}          % dimensionless moment of inertia
\newcommand{\lmax}{\lambda_{\max}}

\title{\bfseries Chaos, its route, and displacement sensitivity\\
in the quasi-steady falling-plate model}

\author{M.\ Subhan Rao\\[2pt]
\small Independent researcher\\[1pt]
\small ORCID: \href{https://orcid.org/0009-0004-4547-141X}{0009-0004-4547-141X}}
\date{\today}

\begin{document}
\maketitle

\begin{abstract}
\noindent
We study the two-dimensional quasi-steady ``falling-card'' model of Andersen, Pesavento \& Wang
(2005) in its thin-card limit, treating the dimensionless moment of inertia $\Is$ as a single control
parameter and asking whether, where, and by what route the model becomes chaotic, and how sensitive
the horizontal landing position is to release conditions. Using a validated numerical pipeline---a
variational (tangent-linear) Lyapunov solver cross-checked against Benettin's method and the
Gottwald--Melbourne 0--1 test, all verified first on the logistic map and the Lorenz system---we
confirm a chaotic band at $\Is \approx [1.63, 2.83]$, with maximal Lyapunov exponent $\lmax$ peaking
at $+0.140$ near $\Is=2.2$, consistent with the value $0.13\pm0.01$ reported in the source paper. The
route into chaos is \emph{not} a Feigenbaum period-doubling cascade: only two successive doublings
occur before the period-two orbit gives way, through a narrow band threaded with periodic windows, to
chaos; power spectra rule out a quasiperiodic (torus) route. Within the chaotic band the standard
deviation of horizontal landing displacement across an ensemble of perturbed release conditions grows
exponentially, at a rate equal to $\lmax$ within estimation uncertainty; outside the band it is
insensitive. Up/down continuation and basin sampling reveal no hysteresis and no non-trivial
multistability. We report one discrepancy with the source paper's figure: at $\Is=1.45$ this model
tumbles with period one, not period two. All findings are properties of the phenomenological model,
not of physical falling objects.

\medskip
\noindent\textbf{Keywords:} falling plate; fluttering; tumbling; deterministic chaos; Lyapunov
exponent; 0--1 test; route to chaos; sensitive dependence.
\end{abstract}

\section{Introduction}\label{sec:intro}

Freely falling plates, cards, leaves, and disks display a well-known sequence of descent styles---
steady descent, side-to-side fluttering, end-over-end tumbling, and irregular motion---that has been
documented for disks~\cite{willmarth64,field97}, strips and cards~\cite{belmonte98,mahadevan99}, and
in direct numerical simulation~\cite{pw04}. Field et al.~\cite{field97} demonstrated genuinely
chaotic dynamics of falling disks, and Belmonte et al.~\cite{belmonte98} mapped the flutter-to-tumble
transition through a Froude-like control parameter. Andersen, Pesavento \& Wang (hereafter
APW)~\cite{apw91,apw65} distilled the two-dimensional dynamics into a low-dimensional quasi-steady
ordinary-differential-equation model in which the fluid forces enter through translational and
rotational drag and a circulatory (``lift'') term, with coefficients fixed from direct force
measurements~\cite{wbd04}. As the dimensionless moment of inertia $\Is$ is varied, the model
reproduces the observed flutter--tumble sequence, and APW reported a chaotic regime with a positive
Lyapunov exponent.

This paper revisits the model as a dynamical system in its own right. Our aims are (i)~to confirm,
with independent diagnostics and pre-verified code, whether a chaotic regime exists and to delimit it;
(ii)~to characterise the \emph{route} by which chaos appears as $\Is$ increases, testing---rather than
assuming---the hypothesis of Feigenbaum period-doubling scaling~\cite{feigenbaum78}; and (iii)~to
quantify the sensitivity of the horizontal landing position to release conditions and connect it to
the Lyapunov exponent. Throughout we treat null and negative results as first-class outcomes: a chaos
claim requires a positive Lyapunov exponent \emph{and} an independent confirmation, and a Feigenbaum
analysis is undertaken only if a genuine doubling cascade is resolved.

\paragraph{Contributions.} Relative to the original model study, this paper (i)~independently confirms
the chaotic regime with two methods---a variational Lyapunov exponent and the 0--1 test---and
delimits it to $\Is\in[1.63,2.83]$; (ii)~resolves the route to chaos, showing it is \emph{not} a
Feigenbaum cascade (only two successive doublings occur) and \emph{not} quasiperiodic; (iii)~
establishes a quantitative link between the Lyapunov exponent and the exponential growth of
landing-position sensitivity; (iv)~reports the absence of hysteresis and of non-trivial
multistability; and (v)~documents a period discrepancy at $\Is=1.45$. All diagnostics are verified on
standard test systems beforehand, and every result is reproducible from a single command.

All findings are properties of this phenomenological model. Its coefficients were tuned by APW for
qualitatively correct behaviour and are not fitted to any specific object; the quasi-steady, purely
two-dimensional closure has a limited range of validity. We therefore make no claim about physical
falling bodies. The complete code, configurations, cached data, and the sources of this manuscript
are openly available (Section~\ref{sec:availability}), and every figure and number below is
regenerated by a single command.

\section{Model}\label{sec:model}

We use the dimensionless thin-card limit ($\beta=b/a\to0$, with $a,b$ the semi-axes) of APW,
equations (5.1)--(5.3) with closures (4.7)--(4.9)~\cite{apw91}. Length is scaled by the semi-major
axis $a$, velocity by $U=\sqrt{(\rho_s/\rho_f-1)\,g\,b}$, and time by $a/U$; here $\rho_s,\rho_f$ are
the solid and fluid densities and $g$ is gravity. The control parameter is $\Is=\rho_s b/(\rho_f a)$.
The reduced state is $\mathbf{u}=(v_x',v_y',\theta,\omega)$, where $v_x',v_y'$ are the centre-of-mass
velocity components in the body-fixed (co-rotating) frame, $\theta$ is the card angle (the only cyclic
coordinate, period $2\pi$), and $\omega=\dot\theta$. The lab-frame position $(x,y)$ follows by
quadrature and does not feed back into the dynamics (translation invariance),
\begin{equation}
\dot x = v_x'\cos\theta - v_y'\sin\theta,\qquad
\dot y = v_x'\sin\theta + v_y'\cos\theta.
\end{equation}
The governing equations are
\begin{align}
\Is\,\dot v_x' &= (\Is{+}1)\,\omega v_y' - \Gamma v_y' - \sin\theta - F_x', \\
(\Is{+}1)\,\dot v_y' &= -\Is\,\omega v_x' + \Gamma v_x' - \cos\theta - F_y', \\
\tfrac14(\Is{+}\tfrac12)\,\dot\omega &= -v_x' v_y' - \tau, \qquad \dot\theta = \omega,
\end{align}
with closures (writing $s=\sqrt{v_x'^2+v_y'^2}$)
\begin{equation}
\Gamma = \frac{2}{\pi}\!\left[-C_T\,\frac{v_x' v_y'}{s} + C_R\,\omega\right],\quad
(F_x',F_y') = \frac{1}{\pi}\!\left[A - B\,\frac{v_x'^2-v_y'^2}{s^2}\right] s\,(v_x',v_y'),\quad
\tau = (\mu_1+\mu_2|\omega|)\,\omega.
\end{equation}
Here $\Gamma$ is the circulation (rotational-lift plus translational term), $\mathbf{F}=(F_x',F_y')$
the fluid drag/lift force, and $\tau$ the dissipative aerodynamic torque. The coefficients, fixed
throughout, are $C_T=1.2$, $C_R=\pi$, $A=1.4$, $B=1.0$, and $\mu_1=\mu_2=0.2$~\cite{apw91,wbd04}. The
direction of $(v_x',v_y')$ is undefined at zero speed; $\Gamma v$ and $\mathbf{F}$ are degree-two
homogeneous and tend to zero there ($C^1$ but not $C^2$), so admissible initial conditions carry
non-zero speed. The torque term $|\omega|\omega$ is likewise $C^1$ but not $C^2$ at $\omega=0$.

The model has two families of fixed points: \emph{edge-on} $(\mp V,0,\pi/2\text{ or }3\pi/2,0)$ with
$V=\sqrt{\pi/(A-B)}\approx2.8025$, and \emph{broadside} $(0,\mp W,0\text{ or }\pi,0)$ with
$W=\sqrt{\pi/(A+B)}\approx1.1441$. The edge-on points are saddles for all $\Is$; the broadside points
lose stability through a Hopf bifurcation, and a heteroclinic connection at $\Is_C\approx1.2191$
mediates the flutter-to-tumble transition.

\section{Numerical methods}\label{sec:methods}

\paragraph{Integration.} Trajectories are integrated with \texttt{scipy.integrate.solve\_ivp} using
the explicit DOP853 method (adaptive high-order Runge--Kutta) for production runs and RK45 for coarse
scans. Production tolerances are $\mathrm{rtol}=10^{-10}$, $\mathrm{atol}=10^{-12}$, fixed by a
tolerance-convergence study (Section~\ref{sec:validation}). Event detection uses the dense output for
sub-step localisation.

\paragraph{Lyapunov exponents.} The maximal exponent $\lmax$ is computed by the variational
(tangent-linear) method: the state and a tangent vector are advanced as one augmented system sharing
identical adaptive steps, with periodic QR reorthonormalisation; $\lmax$ is the time-averaged log
growth rate. As an independent cross-check we use Benettin's two-trajectory renormalisation
method~\cite{benettin80}, implemented as a single augmented reference$+$perturbed system (shared
steps) so that it does not reuse the analytic Jacobian. The uncertainty on $\lmax$ is estimated from
the spread over averaging times.

\paragraph{Independent chaos test.} We apply the Gottwald--Melbourne 0--1 test~\cite{gm04,gm09} to a
scalar observable (the angular velocity sampled on a Poincar\'e section, i.e.\ stroboscopically rather
than at the fine integrator step), reporting the median indicator $K$ over $100$ random frequencies;
$K\to1$ indicates chaos and $K\to0$ regularity.

\paragraph{Poincar\'e section and period counting.} Because a one-signed tumbling orbit never crosses
$\omega=0$, we section on the cyclic angle, recording the reduced state at event-located crossings of
$\theta\equiv\text{const}\,(\mathrm{mod}\,2\pi)$; distinct clustered return points give the period.

\paragraph{Reproducibility.} Every output is content-addressed cached with a provenance sidecar
recording the git hash, Python/NumPy/SciPy versions, parameters, and seed. The cache key includes a
hash of the physics/numerics source files, so any change to the model or solvers invalidates dependent
results. A single command, \texttt{python run\_all.py}, regenerates the entire figure set.

\section{Validation}\label{sec:validation}

No plate result was trusted until a validation ladder passed. \emph{Limiting cases:} with all fluid
coefficients set to zero the model reduces to constant-acceleration free fall along a principal axis,
reproduced to $\sim10^{-9}$; the fixed-point residuals vanish to $<10^{-13}$. \emph{Linearisation:}
the analytic Jacobian agrees with finite differences and, at the heteroclinic value $\Is_C=1.2191$,
gives the edge-on saddle eigenvalues $-0.5854$, $+0.3813$, $-0.9861\pm2.5949\,i$, matching APW; the
decoupled stable eigenvalue equals the closed form $-2\sqrt{(A-B)/\pi}/\Is=-0.58539$.
\emph{Regimes across $\Is_C$:} integration reproduces fluttering at $\Is=1.2190$ and tumbling at
$1.2192$. \emph{Tolerance convergence:} the trajectory is Cauchy-convergent as tolerances tighten, and
DOP853 and RK45 agree at tight tolerance. \emph{Lyapunov code:} verified on the logistic map
($\lambda=\ln2=0.6931$ at $r=4$; the analytic period-two value $-0.9163$ at $r=3.2$) and the Lorenz
system~\cite{lorenz63} ($\lmax=0.907$, a zero exponent, and the exact trace sum
$\sum\lambda_i=-(\sigma+1+\beta)$); the 0--1 test returns $K\approx1$ for chaotic and $K\approx0$ for
periodic logistic series. The full suite comprises $59$ automated tests.

\section{Results}\label{sec:results}

Table~\ref{tab:regimes} summarises the regime sequence and chaos diagnostics; the subsections give the
detail.

\begin{table}[htbp]
\centering
\caption{Regime bands versus $\Is$ with the two chaos diagnostics. The heteroclinic flutter--tumble
transition is at $\Is_C\approx1.2191$; period-doublings occur at $\Is\approx1.474$ and $1.599$.}
\label{tab:regimes}
\begin{tabular}{lccc}
\toprule
$\Is$ range & Regime & $\lmax$ & median $K$ \\
\midrule
$[1.00, 1.20]$ & fluttering & $\approx 0$ & $\approx 0$ \\
$[1.22, 1.47]$ & period-one tumbling & $\approx 0$ & $\approx 0$ \\
$[1.47, 1.60]$ & period-two (flutter/tumble mixture) & $\approx 0$ & $\approx 0$ \\
$[1.63, 2.83]$ & \textbf{chaotic} & up to $+0.140$ & $\approx 1$ \\
$[2.85, 3.00]$ & small-amplitude broadside fluttering & $\approx 0$ & $\approx 0$ \\
\bottomrule
\end{tabular}
\end{table}

\subsection{Regime sequence}\label{sec:e1}

A parallel sweep over $\Is\in[1.0,3.0]$ recovers an ordered sequence of regimes (Fig.~\ref{fig:e1}):
fluttering for $\Is\lesssim1.2$; period-one tumbling from the heteroclinic transition at
$\Is_C\approx1.2191$; a period-two band; an aperiodic (chaotic) band; and a return to small-amplitude
broadside fluttering near $\Is=3.0$. The bifurcation diagram of the section angular velocity displays
the expected structure---paired flutter branches, a single tumbling branch, a period-two split, a
chaotic smear, and flutter again.

\begin{figure}[htbp]
\centering
\includegraphics[width=0.84\linewidth]{e1_regime_map}
\caption{Regime map versus $\Is$. Top: bifurcation diagram of the angular velocity on the $\theta=0$
Poincar\'e section. Bottom: regime strip.}
\label{fig:e1}
\end{figure}

\subsection{Chaos, confirmed by two diagnostics}\label{sec:e2}

The maximal Lyapunov exponent is zero within numerical error ($|\lmax|<10^{-3}$) throughout the
periodic bands and rises to a plateau of $0.12$--$0.14$ across $\Is\approx[1.63,2.83]$, peaking at
$\lmax=+0.140$ at $\Is=2.2$; the independent Benettin estimate there is $0.137$, and the 0--1 test
median $K$ jumps from $0$ to $\approx1$ over the same band (Fig.~\ref{fig:e2}). Requiring \emph{both}
a robustly positive $\lmax$ and $K>0.5$ confirms chaos in $\Is\approx[1.63,2.83]$. The conjunction is
essential: the 0--1 test alone produces spurious elevated $K$ at bifurcation points (transient
artifacts that decay with longer sampling), and finite-time numerical noise produces sub-$10^{-3}$
positive $\lmax$ in periodic bands; each diagnostic vetoes the other's false positives. The peak value
is consistent with APW's reported $0.13\pm0.01$~\cite{apw91}.

\begin{figure}[htbp]
\centering
\includegraphics[width=0.84\linewidth]{e2_lyapunov}
\caption{The chaos gate. Top: bifurcation backdrop. Middle: $\lmax(\Is)$ with uncertainty; the shaded
region marks robustly positive $\lmax$. Bottom: median 0--1 indicator $K(\Is)$ (dashed line at
$K=0.5$).}
\label{fig:e2}
\end{figure}

\subsection{The route to chaos is not Feigenbaum}\label{sec:e3}

Resolving the transition finely (event-located sections, integrator-precision return points) shows
only \emph{two} successive period-doublings: period-1 $\to$ 2 at $\Is\approx1.474$ and 2 $\to$ 4 at
$\Is\approx1.599$ (Fig.~\ref{fig:e3}). There is no period-eight; the period-four state gives way,
across a narrow window ($\sim0.01$ in $\Is$), to chaos interspersed with periodic windows (period-3,
-5, -6). Power spectra distinguish these: the window states show clean line spectra (commensurate
harmonics of a single fundamental) whereas the ``many-branch'' states show a broadband floor identical
to the chaotic control---i.e.\ weak chaos, not two incommensurate frequencies. There is thus no
quasiperiodic (torus / Ruelle--Takens) route. Because a Feigenbaum ratio $\delta$ requires at least
three successive doublings~\cite{feigenbaum78}, and only two are present, we do \emph{not} report a
scaling exponent: doing so from two doublings would be unsupported. The route is best described as a
truncated doubling ($1\to2\to4$) followed by a window-threaded transition to chaos.

\begin{figure}[htbp]
\centering
\includegraphics[width=0.84\linewidth]{e3_route}
\caption{Route to chaos. Top: fine bifurcation zoom on the $\theta=\pi$ section with the two detected
doublings marked (dashed lines). Bottom: $\lmax(\Is)$ over the same window.}
\label{fig:e3}
\end{figure}

\subsection{Landing-displacement sensitivity}\label{sec:e5}

For ensembles of $N=200$ trajectories released from initial conditions perturbed by
$\delta_0=10^{-6}$, we measure the standard deviation $\sigma$ of the horizontal displacement
$\Delta x$ at fixed descent depth $D$ (with a grazing guard on the landing event; Fig.~\ref{fig:e5}).
In the chaotic band $\sigma$ grows exponentially, by six to seven orders of magnitude over $D\le200$
chords; in the regular bands it remains flat within noise, i.e.\ displacement is insensitive there
(these regimes are attracting limit cycles, so perturbations decay---linear growth would require
neutral dynamics).

Measured directly in time, the growth rate $s_t\equiv\mathrm{d}(\log\sigma_x)/\mathrm{d}t$ equals the
Lyapunov exponent to within ${\sim}10\%$ (Table~\ref{tab:e5}; mixed sign), consistent with
$\log\sigma_x\approx\log\sigma_0+\lmax\,t$ to within the combined finite-time uncertainty of the two
estimates (an $N=200$ bootstrap gives slope confidence intervals of width $\sim0.01$--$0.02$, and
$\lmax$ itself carries a few-percent finite-time scatter). The corresponding fixed-depth comparison to
$\lmax/U$, with $U$ the mean descent speed, is looser because it folds in a depth-dependent
time-to-depth conversion. Beyond the exponential phase $\sigma$ continues to grow sub-linearly with
$D$; a clean diffusive $\sqrt{D}$ asymptote is not isolated within the depths studied.

\begin{table}[htbp]
\centering
\caption{Displacement-sensitivity growth in the chaotic band. $\sigma$ growth factor is over
$D\le200$ chords; $s_t=\mathrm{d}(\log\sigma_x)/\mathrm{d}t$ is the exponential-phase rate in time;
$\lmax$ is the variational value at $T=4000$. Regular controls $\Is=1.1,1.4$ show growth factors
$\sim30$ (noise) and are flat.}
\label{tab:e5}
\begin{tabular}{ccccc}
\toprule
$\Is$ & $\sigma$ growth factor & $s_t$ & $\lmax$ & $s_t/\lmax$ \\
\midrule
$2.0$ & $3.3\times10^{6}$ & $0.130$ & $0.117$ & $1.11$ \\
$2.2$ & $4.6\times10^{6}$ & $0.132$ & $0.143$ & $0.92$ \\
$2.5$ & $1.2\times10^{7}$ & $0.115$ & $0.127$ & $0.90$ \\
\bottomrule
\end{tabular}
\end{table}

\begin{figure}[htbp]
\centering
\includegraphics[width=0.96\linewidth]{e5_sensitivity}
\caption{Landing-displacement spread. Left: $\sigma(D)$ (semilog) for chaotic (solid) and regular
(dashed) $\Is$, with exponential-phase fits (dotted). Right: chaotic curves (log-log) against
$\sqrt{D}$ and linear references.}
\label{fig:e5}
\end{figure}

\subsection{No non-trivial multistability}\label{sec:e6}

Up and down continuation sweeps of $\langle|\omega|\rangle$ coincide across $\Is\in[1.0,3.0]$: there
is no hysteresis, including at the flutter--tumble transition (Fig.~\ref{fig:e6}). Basin sampling with
random initial conditions finds a single attractor at every diagnostic $\Is$, inside the periodic
windows and the chaotic band alike. The only multistability is the trivial
clockwise/counter-clockwise reflection-symmetry pair (mirror spin directions with identical
statistics). This rules out the possibility that the chaotic diagnostics reflect coexisting attractors
rather than a single chaotic one.

\begin{figure}[htbp]
\centering
\includegraphics[width=0.84\linewidth]{e6_multistability}
\caption{Multistability probe. Top: up (filled) versus down (open) continuation of
$\langle|\omega|\rangle$ versus $\Is$---the branches coincide (no hysteresis). Bottom: basin sampling
(random initial conditions) at diagnostic $\Is$; a single cluster at each.}
\label{fig:e6}
\end{figure}

\section{Discussion}\label{sec:discussion}

The model possesses a genuine chaotic band, confirmed by independent methods and delimited more
sharply than by regime inspection alone. Two aspects refine or depart from expectations. First, the
\emph{route} is not the period-doubling cascade one might anticipate from the early doubling: a single
robust doubling and a second, closely followed by a window-threaded transition, is a qualitatively
different scenario for which Feigenbaum universality does not apply. Second, we find a clean
operational link between a dynamical invariant and an observable of practical interest---the
exponential rate of landing-position spread equals the Lyapunov exponent---while emphasising that the
regular regimes are \emph{insensitive}, so the predictability of landing position depends
discontinuously on $\Is$.

We note one discrepancy with the source paper. APW's Fig.~3 reads period-two tumbling at $\Is=1.45$;
in our integration this value is period-\emph{one} (the section state and revolution time are constant
to $\sim10^{-12}$ across revolutions, and three distinct initial conditions converge to the same orbit
with no bistability). The first doubling occurs here at $\Is\approx1.474$. We report this as a
property of the model under our conventions rather than reconciling it by parameter adjustment; the
initial conditions used for APW's figures are not specified, and small differences in the bifurcation
location are plausible.

\paragraph{Limitations.} The coefficients are phenomenological and the closure quasi-steady and
two-dimensional; results outside the parameter range where that closure is credible should not be
over-interpreted. The plate Lyapunov exponent is only moderately converged at attainable averaging
times (a few-percent finite-time scatter), which sets the precision of the $\lmax$--sensitivity
comparison. The diffusive tail of $\sigma(D)$ and the fine structure of the $[1.598,1.61]$ transition
(intermittency versus crisis) remain open to deeper characterisation but do not affect the conclusions
above.

\section{Conclusions}\label{sec:conclusions}

For the thin-card APW model with the standard coefficients, as $\Is$ increases the descent passes from
fluttering through tumbling to a chaotic band at $\Is\approx[1.63,2.83]$ ($\lmax$ peak $+0.140$),
confirmed by variational and Benettin Lyapunov exponents and the 0--1 test. The route is a truncated
period-doubling (two doublings) followed by windowed chaos---not a Feigenbaum cascade and not a torus.
Horizontal landing displacement is exponentially sensitive to release conditions in the chaotic band,
at the Lyapunov rate, and insensitive elsewhere. The model exhibits no hysteresis and no non-trivial
multistability. We additionally flag a period discrepancy with the source paper at $\Is=1.45$. Natural
extensions are to resolve the $\sqrt{D}$ diffusive tail at larger fall depths and to classify the
narrow $[1.598,1.61]$ transition zone in detail.

\section*{Data and code availability}\label{sec:availability}
All source code, experiment configurations, cached data with provenance sidecars, and the sources of
this manuscript accompany the paper. \texttt{python run\_all.py} regenerates every figure (from cache
or from scratch); \texttt{make paper} builds this document; and an automated test suite
(\texttt{pytest}, $59$ tests) covers the physics and the diagnostics. The Feigenbaum stage is
intentionally not run and is reported as the null result of Section~\ref{sec:e3}.

\section*{Acknowledgements}
The author conceived and directed this study, made all scientific decisions, verified the validation
anchors, interpreted the results, and prepared this manuscript, and is solely responsible for its
content. Software development and manuscript preparation were assisted by general-purpose AI tools;
the author reviewed all code, results, and text. All results are reproducible from the accompanying
code.

\section*{Competing interests}
The author declares no competing interests.

\section*{Funding}
This research received no external funding.

\begin{thebibliography}{99}
\bibitem{willmarth64} W.~W.~Willmarth, N.~E.~Hawk, R.~L.~Harvey,
\emph{Steady and unsteady motions and wakes of freely falling disks},
Phys.\ Fluids \textbf{7}, 197--208 (1964).
\bibitem{field97} S.~B.~Field, M.~Klaus, M.~G.~Moore, F.~Nori,
\emph{Chaotic dynamics of falling disks}, Nature \textbf{388}, 252--254 (1997).
\bibitem{belmonte98} A.~Belmonte, H.~Eisenberg, E.~Moses,
\emph{From flutter to tumble: inertial drag and Froude similarity in falling paper},
Phys.\ Rev.\ Lett.\ \textbf{81}, 345--348 (1998).
\bibitem{mahadevan99} L.~Mahadevan, W.~S.~Ryu, A.~D.~T.~Samuel,
\emph{Tumbling cards}, Phys.\ Fluids \textbf{11}, 1--3 (1999).
\bibitem{pw04} U.~Pesavento, Z.~J.~Wang,
\emph{Falling paper: Navier--Stokes solutions, model of fluid forces, and center of mass elevation},
Phys.\ Rev.\ Lett.\ \textbf{93}, 144501 (2004).
\bibitem{apw91} A.~Andersen, U.~Pesavento, Z.~J.~Wang,
\emph{Analysis of transitions between fluttering, tumbling and steady descent of falling cards},
J.\ Fluid Mech.\ \textbf{541}, 91--104 (2005).
\bibitem{apw65} A.~Andersen, U.~Pesavento, Z.~J.~Wang,
\emph{Unsteady aerodynamics of fluttering and tumbling plates},
J.\ Fluid Mech.\ \textbf{541}, 65--90 (2005).
\bibitem{wbd04} Z.~J.~Wang, J.~M.~Birch, M.~H.~Dickinson,
\emph{Unsteady forces and flows in low Reynolds number hovering flight: two-dimensional
computations vs robotic wing experiments}, J.\ Exp.\ Biol.\ \textbf{207}, 449--460 (2004).
\bibitem{gm04} G.~A.~Gottwald, I.~Melbourne,
\emph{A new test for chaos in deterministic systems},
Proc.\ R.\ Soc.\ Lond.\ A \textbf{460}, 603--611 (2004).
\bibitem{gm09} G.~A.~Gottwald, I.~Melbourne,
\emph{On the implementation of the 0--1 test for chaos},
SIAM J.\ Appl.\ Dyn.\ Syst.\ \textbf{8}, 129--145 (2009).
\bibitem{benettin80} G.~Benettin, L.~Galgani, A.~Giorgilli, J.-M.~Strelcyn,
\emph{Lyapunov characteristic exponents for smooth dynamical systems and for Hamiltonian systems; a
method for computing all of them}, Meccanica \textbf{15}, 9--30 (1980).
\bibitem{lorenz63} E.~N.~Lorenz, \emph{Deterministic nonperiodic flow},
J.\ Atmos.\ Sci.\ \textbf{20}, 130--141 (1963).
\bibitem{feigenbaum78} M.~J.~Feigenbaum,
\emph{Quantitative universality for a class of nonlinear transformations},
J.\ Stat.\ Phys.\ \textbf{19}, 25--52 (1978).
\end{thebibliography}

\end{document}
